% GENERATED FILE. Edit canonical lessons, then run npm run generate:books.
% canonical-source-hash: fnv1a64:b8472dab790e5d6d
% canonical-lessons: HI-C36-ghar, HI-C36-darvaaza, HI-C36-kamra, HI-C36-kursi

\chapter{One House, Four Languages}
\label{ch:hi-one-house-four-languages}
\hlchaptermodality{36}

\begin{chapteropening}
\textbf{By the end of this chapter:} \emph{I can name a house, its door, a room and the chair I am offered, give each one the possessive its gender demands, and say where each of the four words came from — Sanskrit, Persian, Portuguese and Arabic.}

Run \hi{घर}, \hi{दरवाज़ा}, \hi{कमरा} and \hi{कुर्सी} with the right possessive — and hear \hi{मेरा} turn into \hi{मेरी} at the first feminine noun, the promise chapter 2 made and never kept. Then name the four languages one room drew on, spell \hi{कुर्सी} with its virāma and its conjunct, and say \hi{नमस्ते} at the door coming in and \hi{शुभ} \hi{रात्रि} going out.
\end{chapteropening}

\section[\hi{घर}]{\hi{घर} (ghar) — \textquotedblleft{}house,\textquotedblright{} and the fence that came first}

\label{lesson:HI-C36-ghar}



\begin{quote}
You can say your name. Now say where you live — and meet the one thing every Hindi noun carries that English nouns do not.
\end{quote}

\subsection*{The letters in this word}
Two letters. \textbf{\hi{घ}} (\emph{gha}), the breathy pair of \emph{ga}, then \textbf{\hi{र}} (\emph{ra}), already yours from \textbf{\hi{मेरा}}.

\textbf{\hi{र}} is the spineless one: no upright stroke of its own, so the \textbf{\hi{शिरोरेखा}}, the head-line, is all that ties it to its neighbour. Both letters keep their inherent \emph{a} $\to$ \textbf{\hi{घर}}, \emph{ghar}.

\begin{grammarlens}[title={Grammar Lens: the word arrives with a gender}]
\textbf{\hi{घर}} (\emph{ghar}) = \textquotedblleft{}\textbf{house, home}.\textquotedblright{} And it is \textbf{masculine}.

Every Hindi noun is masculine or feminine. Not the thing — the \textbf{word}. There is nothing manly about a house; the gender is a fact about \emph{ghar} the way its spelling is, and you learn the two together or you learn them twice.

It is not decoration: it changes the words around the noun. Because \emph{ghar} is masculine, \textquotedblleft{}my house\textquotedblright{} is \textbf{\hi{मेरा} \hi{घर}} (\emph{merā ghar}), the masculine \textbf{\hi{मेरा}}, exactly as in \textbf{\hi{मेरा} \hi{नाम}}.

Whole sentence: \textbf{\hi{यह} \hi{मेरा} \hi{घर} \hi{है}} (\emph{yah merā ghar hai}), \textquotedblleft{}\textbf{this is my house}.\textquotedblright{}
\end{grammarlens}

\begin{cousinweb}
\textbf{\hi{घर}} comes from Sanskrit \textbf{\hi{गृह}} (\emph{gṛha}), worn smooth through Prakrit \emph{ghara} into the short word you just said. Hindi kept \textbf{\hi{गृह}} too, borrowed back out of Sanskrit for formal use, so both shapes survive.

Further back, \emph{gṛha} descends from an ancient root meaning \textbf{enclosure}, \emph{\textbf{g\textsuperscript{h}erd\textsuperscript{h}-}} — and English holds two plain descendants of it: \textbf{yard} and \textbf{garden}. Slavic holds the \emph{-grad} in \textbf{Belgrade}, \textquotedblleft{}town.\textquotedblright{}

So the oldest sense is not walls and a roof. It is \textbf{a fenced-in piece of ground}. You enclose the land first; the house is what goes inside afterwards.
\end{cousinweb}

\subsection*{Your turn}
Say these aloud:

\begin{itemize}
\raggedright
  \item \emph{ghar} — house; then \emph{merā ghar}; then \emph{yah merā ghar hai}
  \item \emph{merā nām} and \emph{merā ghar} one after the other — the same \emph{merā}, because both nouns are masculine
  \item the two letters, \emph{gha} and \emph{ra}, and which of them has no spine
  \item \emph{gṛha}, then \emph{ghar} — the Sanskrit shape and the worn one; then the old meaning, \emph{enclosure}
\end{itemize}

\begin{tcolorbox}[breakable,colback=teal!4,colframe=teal!35!black,title={Before you move on}]
Is \textbf{\hi{घर}} masculine or feminine, and how do you know from the sentence? (\textbf{Masculine} — \emph{merā}, not \emph{merī}.) What did the root behind \emph{gṛha} mean before it meant a building? (\textbf{An enclosure} — the root behind \emph{yard} and \emph{garden}.) Which of \textbf{\hi{घ}} and \textbf{\hi{र}} hangs from the head-line without a spine of its own? (\textbf{\hi{र}}.)
\end{tcolorbox}
\section[\hi{दरवाज़ा}]{\hi{दरवाज़ा} (darvāzā) — the door, and the oldest word you already know}

\label{lesson:HI-C36-darvaaza}



\begin{quote}
You have a house. This is what you stand at to be let into one — and it is where \textbf{\hi{नमस्ते}} is said and where \textbf{\hi{शुभ} \hi{रात्रि}} is said back.
\end{quote}

\subsection*{The letters in this word}
\textbf{\hi{द}·\hi{र}·\hi{वा}·\hi{ज़ा}} — and the \textbf{\hi{ा}} you already write does two of the four.

The new mark is the dot under \textbf{\hi{ज़}}: the \textbf{nuqta}. Plain \textbf{\hi{ज}} is \emph{ja}. Persian arrived carrying a \textbf{z} Devanagari had no letter for, so the script put a dot beneath the nearest shape and made one — \textbf{\hi{ज़}}, \emph{za}. The same dot makes \textbf{\hi{फ़}} (\emph{fa}) and \textbf{\hi{ख़}}, the sound in \textbf{\hi{ख़ुशी}}.

A nuqta is a small confession in ink: \emph{this sound came from somewhere else.}

\begin{grammarlens}[title={Grammar Lens: the shape that usually tells you}]
\textbf{\hi{दरवाज़ा}} (\emph{darvāzā}) = \textquotedblleft{}\textbf{door},\textquotedblright{} and it is \textbf{masculine} — so \textbf{\hi{मेरा} \hi{दरवाज़ा}}.

Here is your first piece of help. A noun ending in long \textbf{-ā} is usually masculine, and \emph{darvāzā} ends in \textbf{-\hi{ा}}. Keep it as a lean, not a law: later chapters will show you where it fails, in both directions.

\textbf{\hi{घर}} (\emph{ghar}) ends in a consonant and is masculine too — so a consonant ending tells you nothing at all. That one you simply learn.
\end{grammarlens}

\begin{cousinweb}
\textbf{\hi{दरवाज़ा}} is \textbf{Persian}, \textbf{\ar{دروازه}} (\emph{darvāza}), and its first syllable is where this gets remarkable. Persian \textbf{\ar{در}} (\emph{dar}) means \textquotedblleft{}\textbf{door},\textquotedblright{} and it descends from the ancient root \emph{\textbf{d\textsuperscript{h}wer-}} — as do English \textbf{door}, Latin \textbf{forēs}, Greek \textbf{thúrā} and Sanskrit \textbf{\hi{द्वार}} (\emph{dvāra}).

Now the twist. Hindi \emph{also} has \textbf{\hi{द्वार}} (\emph{dvār}), taken straight out of Sanskrit for formal use. So the borrowed Persian word and the learned Sanskrit one are \textbf{the same ancient word, arrived twice}, by two roads, centuries apart.

The \textbf{-vāza} part is less settled: usually connected to a Persian word for \textquotedblleft{}\textbf{open},\textquotedblright{} but disputed, and worth leaving open.
\end{cousinweb}

\subsection*{Your turn}
Say these aloud:

\begin{itemize}
\raggedright
  \item \emph{darvāzā}, then \emph{merā darvāzā}, then \emph{merā ghar} — one \emph{merā}, two masculine nouns
  \item the arrival and the leaving at the same door — \emph{namaste} going in, \emph{shubh rātri} going out
  \item \emph{ja}, then \emph{za} — the letter, then the letter with the dot
  \item \emph{dar}, \emph{door}, \emph{thúrā}, \emph{dvāra} — one word, four languages
\end{itemize}

\begin{tcolorbox}[breakable,colback=teal!4,colframe=teal!35!black,title={Before you move on}]
What does the dot under \textbf{\hi{ज़}} tell you? (That the sound is \textbf{borrowed} — Devanagari had no \emph{z}.) Persian \emph{dar} and English \emph{door} — coincidence or cousins? (\textbf{Cousins}, from \emph{\textbf{d\textsuperscript{h}wer-}} — and Hindi's own \textbf{\hi{द्वार}} is the third.) Which ending gives you a fair guess at masculine? (\textbf{Long -ā}, as in \emph{darvāzā}. A consonant ending, as in \emph{ghar}, tells you nothing.)
\end{tcolorbox}
\section[\hi{कमरा}]{\hi{कमरा} (kamrā) — the room that is also a camera}

\label{lesson:HI-C36-kamra}



\begin{quote}
Through the door, and into this. A short, ordinary word — which travelled to India by sea, on Portuguese ships.
\end{quote}

\subsection*{The letters in this word}
\textbf{\hi{क}·\hi{म}·\hi{रा}} — and \textbf{\hi{म}} (\emph{ma}) is one of the first two letters you ever wrote, beside \textbf{\hi{न}}.

Both of those hang from the head-line and both stand on a right-hand spine; that upright is what your eye follows along a line of Devanagari. \textbf{\hi{क}} and \textbf{\hi{र}} you have as well, so this word costs you nothing new — only the \textbf{\hi{ा}} riding on the final \textbf{\hi{र}}.

\begin{grammarlens}[title={Grammar Lens: the lean holds}]
\textbf{\hi{कमरा}} (\emph{kamrā}) = \textquotedblleft{}\textbf{room},\textquotedblright{} and it is \textbf{masculine} — \textbf{\hi{मेरा} \hi{कमरा}}.

Long \textbf{-ā}, masculine again, exactly as \textbf{\hi{दरवाज़ा}} was. Three nouns in, and the possessive has not changed once: \emph{merā ghar}, \emph{merā darvāzā}, \emph{merā kamrā}.

Hold on to that flatness. It is about to break.
\end{grammarlens}

\begin{cousinweb}
\textbf{\hi{कमरा}} is neither Sanskrit nor Persian. It is \textbf{Portuguese} — \textbf{câmara}, \textquotedblleft{}chamber, room\textquotedblright{} — left behind on India's western coast by the traders who ran that shore from the sixteenth century, along with \textbf{\hi{चाबी}} (\emph{chābī}, \textquotedblleft{}key\textquotedblright{}) and \textbf{\hi{अलमारी}} (\emph{almārī}, \textquotedblleft{}cupboard\textquotedblright{}).

Portuguese \emph{câmara} is Latin \textbf{camera}, \textquotedblleft{}a vaulted chamber,\textquotedblright{} and Latin took it from Greek \textbf{kamára} — where the trail honestly stops; the Greek word's own origin is unknown.

Then the good part. English \textbf{camera} is that same Latin word: a \emph{camera obscura} was literally a \textquotedblleft{}dark room,\textquotedblright{} and when the dark room shrank into a box the name went with it. And Hindi later borrowed it back through English as \textbf{\hi{कैमरा}} (\emph{kaimrā}) for the device.

So \textbf{\hi{कमरा}} and \textbf{\hi{कैमरा}} are the \emph{same Latin word}, arrived twice by two different routes — the second time carrying a machine that did not exist when the first one landed.
\end{cousinweb}

\subsection*{Your turn}
Say these aloud:

\begin{itemize}
\raggedright
  \item \emph{kamrā}, then \emph{merā kamrā}
  \item the three so far, in one breath — \emph{merā ghar, merā darvāzā, merā kamrā}
  \item \emph{kamrā} and \emph{kaimrā} — one Latin word, two arrivals
  \item where each of the three came from — Sanskrit, Persian, Portuguese
\end{itemize}

\begin{tcolorbox}[breakable,colback=teal!4,colframe=teal!35!black,title={Before you move on}]
Which language handed Hindi \textbf{\hi{कमरा}}, and how did it get there? (\textbf{Portuguese}, by sea, on the western coast.) What does \emph{kamrā} share with English \emph{camera}? (\textbf{One Latin word}, \emph{camera}, \textquotedblleft{}vaulted chamber\textquotedblright{} — Hindi has it twice, as \emph{kamrā} and \emph{kaimrā}.) Say \textquotedblleft{}my house, my door, my room.\textquotedblright{} (\emph{Merā ghar, merā darvāzā, merā kamrā}.)
\end{tcolorbox}
\section[kursī]{\hi{कुर्सी} (kursī) — sit down, and watch \hi{मेरा} change}

\label{lesson:HI-C36-kursi}



\begin{quote}
House, door, room — and now what you are offered once inside. This word breaks the pattern the other three set.
\end{quote}

\subsection*{The letters in this word}
\textbf{\hi{कु}·\hi{र्}·\hi{सी}} — three old letters, \textbf{\hi{क}}, \textbf{\hi{र}}, \textbf{\hi{स}}, and two working marks.

\textbf{\hi{क}} takes \textbf{\hi{ु}} for \emph{ku}. Then \textbf{\hi{र}} with the \textbf{virāma}, the stroke that kills a consonant's built-in \emph{a}: a bare \textbf{\hi{र्}} cannot stand alone, so it fuses with the \textbf{\hi{स}}, and the hook above \textbf{\hi{स}} \emph{is} that \textbf{\hi{र}} — the fusing you met in \textbf{\hi{स्त}}.

Then \textbf{\hi{सी}}: the long \textbf{\hi{ी}} stands after its consonant and sounds after it, while \textbf{\hi{ि}} stands before and still sounds after. One of the two lies to your eye.

\begin{grammarlens}[title={Grammar Lens: your first feminine noun}]
\textbf{\hi{कुर्सी}} (\emph{kursī}) = \textquotedblleft{}\textbf{chair},\textquotedblright{} and it is \textbf{feminine}.

So \textbf{\hi{मेरी} \hi{कुर्सी}}, not \emph{merā}. When you learnt \textbf{\hi{मेरा}} you were told you would meet \textbf{\hi{मेरी}} at your first feminine noun. Here it is.

Nothing about a chair is female. The gender belongs to the \textbf{word}, and it changes \textbf{\hi{मेरा}}. Get it wrong and every word beside the noun goes wrong too.
\end{grammarlens}

\begin{grammarlens}[title={Grammar Lens: the lean, and where it leans wrong}]
The rule of thumb: \textbf{long -ā tends masculine, long -ī tends feminine, a consonant ending tells you nothing.} A lean, not a law: \emph{darvāzā}, \emph{kamrā} and \emph{kursī} obey it, \emph{ghar} sits outside it, and later words will disobey it outright.
\end{grammarlens}

\begin{cousinweb}
\textbf{\hi{कुर्सी}} is \textbf{Arabic}, \textbf{\ar{كرسي}} (\emph{kursī}), \textquotedblleft{}seat, throne,\textquotedblright{} reaching Hindi through \textbf{Persian}, the road most Arabic words took.

Arabic did not invent it either. It has passed along the Near East for four thousand years: Sumerian \textbf{guza}, Akkadian \textbf{kussû}, then Ugaritic, Hebrew and Aramaic, Arabic most likely taking it from \textbf{Aramaic}. Not one family — the word crossed between unrelated families because the object did.

English \textbf{chair}? Greek \textbf{kathédra}, through Latin and French. No relation at all.

And what this chapter has done: \textbf{\hi{घर}} Sanskrit, \textbf{\hi{दरवाज़ा}} Persian, \textbf{\hi{कमरा}} Portuguese, \textbf{\hi{कुर्सी}} Arabic. One room, four languages — that is Hindi, not an accident of these four words.
\end{cousinweb}

\subsection*{Your turn}
Say these aloud:

\begin{itemize}
\raggedright
  \item \emph{merī kursī}, then \emph{merā kamrā} — one possessive, two shapes
  \item the chapter — \emph{merā ghar, merā darvāzā, merā kamrā, merī kursī}
  \item each word's source — Sanskrit, Persian, Portuguese, Arabic
  \item \emph{namaste} at the \emph{darvāzā}, then \emph{shubh rātri} at the same door
\end{itemize}

\begin{tcolorbox}[breakable,colback=teal!4,colframe=teal!35!black,title={Before you move on}]
Why \textbf{\hi{मेरी} \hi{कुर्सी}}? (\textbf{\hi{कुर्सी} is feminine}; the possessive agrees with the noun.) How far back does \emph{kursī} go, and did it stay in one family? (\textbf{Sumerian} — and no, it crossed unrelated families.) Name the chapter's four sources. (\textbf{Sanskrit, Persian, Portuguese, Arabic}.)
\end{tcolorbox}
